% !TEX root = ../main.tex
\section{Methods}
\label{sec:methods}

\subsection{Inventory binning}
\label{sec:methods_binning}

The workflow operates on fixed-area permanent sample plots from the Qu\'ebec provincial inventory. For each species-group / cover-type combination we aggregate tallies into 2~cm classes spanning 10--60~cm DBH (diameter at breast height). Each 0.04~ha plot receives the constant expansion factor $s_{\text{plot}} = 25$, yielding a stand table $\hat{y}_i = s_{\text{plot}} t_i$ for bin $i$, where $t_i$ is the mean tally across plots. We then normalise the stand table so that $\sum_{i \in I} W \hat{y}_i^{(n)} = 1$ for bin width $W = 2$~cm, producing a truncated empirical density on $[a, b] = [10, 60]$~cm.
The upper bound $b$ reflects the observed maximum in the PSP tallies (and management caps where applicable); if no defensible upper limit exists, the workflow can be applied with a larger $b$ or with left truncation only.

\subsection{Weighted non-linear least squares}
\label{sec:methods_nlls}

We fit parametric distributions by matching the empirical truncated density with candidate PDFs at the bin midpoints using weighted non-linear least squares. The objective is
\begin{equation}
  Z = \min_{\bm{P}} \sum_{i \in I} e_i^2,
\end{equation}
with residuals
\begin{equation}
  e_i = w_i \left[f(x_i; \bm{P}) - \hat{y}_i^{(n)}\right],
\end{equation}
where $x_i$ is the class midpoint, $f$ the density evaluated at $x_i$, and $\bm{P}$ the free parameters. We retain unit bin weights ($w_i = 1$) to streamline the narrative and align the comparison with the published truncated baseline. The companion repository notebook documents the plot-level tallies and demonstrates how to construct inverse-variance weights when analysts wish to down-weight sparse edge bins.
Optimisation uses the Levenberg--Marquardt algorithm as implemented in \texttt{lmfit} \citep{levenberg1944method, marquardt1963algorithm}.

Truncated datasets exhibit biased residual sums when complete-form densities are fit directly. The recommended approach rescales the target density to produce the truncated form
\begin{equation}
  f^{T}(x; \bm{P}) =
    \begin{cases}
      c\, f(x; \bm{P}) & a \leq x \leq b, \\
      0 & \text{otherwise},
    \end{cases}
\end{equation}
with scaling constant $c = [F(b; \bm{P}) - F(a; \bm{P})]^{-1}$ and $F$ the cumulative distribution function. When a closed form for $F$ is not available we evaluate the integral numerically.

\subsection{Two-stage estimator}
\label{sec:methods_twostage}

The two-stage estimator includes a positive scaling parameter $s$ so the complete-form fit honours the truncated support. Stage one solves
\begin{equation}
  \min_{\bm{P},\, s>0} \sum_{i \in I} \left[s f(x_i; \bm{P}) - \hat{y}_i^{(n)}\right]^2,
\end{equation}
returning provisional parameters $\bm{P}^{(1)}$ and $s^{(1)}$. Stage two holds $s^{(1)}$ fixed, refits $\bm{P}$ with the same objective, and yields the final estimates $\bm{P}^{(2)}$.
Because $f$ integrates to one, the fitted scaling parameter provides an implied retained mass on $[a,b]$ via $F(b;\bm{P})-F(a;\bm{P}) \approx 1/s$; in the limit of a perfect match to the truncated form, $s$ equals the normalising constant $c$.

\subsection{Comparative fitting workflow}
\label{sec:methods_benchmark}

We compare three fitting strategies for each meta-plot:
\begin{enumerate}
  \item \textbf{1sc} -- a complete-form fit with the scaling fixed to $s = 1$;
  \item \textbf{1st} -- a truncated likelihood fit using $f^{T}(x; \bm{P})$ on $[a, b]$;
  \item \textbf{2sc} -- the two-stage complete-form procedure described above.
\end{enumerate}
The 1sc configuration is retained intentionally as an operational baseline: in many applied inventory workflows, analysts still fit complete-form densities directly to truncated grouped data with no explicit truncation correction. Including 1sc therefore quantifies the practical consequence of omitting truncation handling, while 1st serves as the reference truncated fit and 2sc as the proposed complete-form alternative.
We wrap each probability density function into an \texttt{lmfit} \texttt{Model}, initialise parameters from empirical moments, and minimise $Z$ for all three configurations.

\subsection{Candidate distributions and model selection}
\label{sec:methods_candidates}

We concentrate on the two-parameter Weibull and gamma distributions, both specialisations of the generalised gamma family that cover common stand structures with few parameters. Each species-group / cover-type meta-plot is fit with both distributions under the 1sc, 1st, and 2sc regimes. We compute the small-sample Akaike information criterion,
\begin{equation}
  \text{AICc} = \text{AIC} + \frac{2K(K+1)}{n - K - 1},
\end{equation}
where $n$ is the number of non-zero bins and $K$ counts the number of free parameters plus the implicit error variance term. We adjust $K$ to account for the additional scaling terms associated with each configuration: $K_{1sc} = |\bm{P}| + 1$, $K_{1st} = |\bm{P}| + 3$ (two truncation bounds plus variance), and $K_{2sc} = |\bm{P}| + 2$. The best model for each meta-plot is selected by minimising AICc. We report parameter estimates, associated standard errors, and AICc values for both stage-one and stage-two fits to facilitate comparison with the truncated baseline.

\subsection{Robustness assessment workflow}
\label{sec:methods_robustness}

The robustness assessment has two components. First, all 32 species-group / cover-type combinations are fit and empirical histograms are classified into shape classes using a deterministic peak-location rule after light smoothing: inverse-J-like (peak in the first 20\% of bins), mid-unimodal (20--60\%), and right-peaked/flat ($>$60\%). Method win counts (AICc), 1st--2sc discrepancy metrics, and fit-failure counts are then summarised by shape class.

Second, a canonical synthetic benchmark is evaluated with four scenario families (inverse-J, mid-unimodal, near-flat, and right-shifted) under fixed random seeds and replicated sample-size regimes. For each synthetic dataset, the same 1sc/1st/2sc workflow is run and win rates and discrepancy metrics are aggregated.
Compact robustness summaries are provided in Appendix Tables~\ref{tab:app_shape_wins}--\ref{tab:app_sim_summary}.
